% Generated draft: requires mathematical and source review.
\documentclass[11pt]{article}
\usepackage[margin=25mm]{geometry}
\usepackage{fontspec}
\setmainfont{Latin Modern Roman}
\usepackage{amsmath,amssymb,mathtools,unicode-math}
\setmathfont{Latin Modern Math}
\usepackage{enumitem,needspace,xurl,hyperref}
\setlength{\parindent}{0pt}
\setlength{\parskip}{5pt}
\setlength{\emergencystretch}{3em}
\newcommand{\N}{\mathbb{N}}
\newcommand{\Z}{\mathbb{Z}}
\newcommand{\Q}{\mathbb{Q}}
\newcommand{\R}{\mathbb{R}}
\newcommand{\C}{\mathbb{C}}
\newcommand{\F}{\mathbb{F}}
\newcommand{\A}{\mathbb{A}}
\newcommand{\PP}{\mathbb{P}}
\newcommand{\E}{\mathbb{E}}
\newcommand{\eps}{\varepsilon}
\newcommand{\dd}{\,\mathrm{d}}
\begin{document}
\section*{UnsolvedMath ID: 30005217}
\textbf{Source problem code:} catalogue entry
\subsection{Definitions and mathematical statement}
The supplied record presents four broad research questions rather than a conjecture with specified hypotheses and quantifiers. Its title refers to finite group characters, but the quoted statement does not specify a class of groups, representations, or characters. The following definitions provide a conditional finite-group interpretation; they do not resolve the missing scope.

Let \(G\) be a finite group. For a finite-dimensional complex representation \(V\) of \(G\), its character is the function
\[
\chi_V\colon G\longrightarrow\C,\qquad \chi_V(g)=\operatorname{tr}(g\mid V).
\]
Write \(\operatorname{Irr}_{\C}(G)\) for the set of irreducible complex characters of \(G\). If \(V\) and \(W\) are complex representations, give \(V\otimes_{\C}W\) the diagonal action \(g(v\otimes w)=(gv)\otimes(gw)\). Its character is the pointwise product \(\chi_V\chi_W\). Denote its irreducible multiplicities by the nonnegative integers \(m_{\chi_V,\chi_W}^{\psi}\) in
\[
\chi_V\chi_W=\sum_{\psi\in\operatorname{Irr}_{\C}(G)}m_{\chi_V,\chi_W}^{\psi}\,\psi.
\]
For a complex character \(\chi\), define its vanishing set by
\[
Z_G(\chi)=\{g\in G:\chi(g)=0\}.
\]
For a prime \(p\), a field \(k\) of characteristic \(p\), and a finite-dimensional \(kG\)-module \(M\), its representation dimension is \(\dim_k M\). If ``modular character'' means a Brauer character \(\varphi_M\) in an appropriate splitting-field setting, its domain consists of the \(p\)-regular elements, namely those whose orders are not divisible by \(p\), and its degree is \(\varphi_M(1)=\dim_k M\). The source does not specify this interpretation or these field hypotheses.

With these qualifications, the supplied research questions are:
\begin{enumerate}
\item Describe the decomposition of the tensor product of a pair of complex characters, or equivalently determine its irreducible multiplicities in the intended setting.
\item Understand the vanishing sets of complex characters.
\item Determine the asymptotic behaviour of complex characters.
\item Determine the dimensions associated with modular characters.
\end{enumerate}
The source leaves unspecified whether the characters are irreducible, which groups or families are intended, what form of description is sought, which parameters tend to a limit, how character values are normalized, and which modular representations and characteristics are under consideration. Consequently, these questions cannot be converted from the supplied evidence into a single fully quantified mathematical assertion. In particular, no asymptotic limit, estimate, or modular dimension formula is asserted here.
\subsection{Short English statement}
The record asks how tensor products of complex characters decompose, where characters vanish, how characters behave asymptotically, and what dimensions occur for modular characters. It supplies a research programme whose group families, limiting regimes, and modular hypotheses remain unspecified. The finite-group definitions above are a conditional interpretation motivated by the supplied title. This is a statement-only draft requiring human mathematical review.
\subsection{Sources}
\begin{itemize}
\item Supplied local UnsolvedMath registry record, permanent ID \(30005217\), titled \emph{Structure and Asymptotics of Finite Group Characters}; exact source problem code \texttt{catalogue entry}. The record belongs to the Oberwolfach Reports Open Problems collection. No individual proposer is supplied; the recorded proposed year is \(2022\).
\item Supplied report citation: \emph{Character Theory and Categorification (2023)}, Oberwolfach Reports, DOI \texttt{10.4171/owr/2022/39}, pp.~2267--2333. Source URL and literature-source link labelled ``original Oberwolfach report'': \url{https://doi.org/10.4171/owr/2022/39}. The supplied background describes Oberwolfach Reports as documenting mathematical workshops and their open-problem discussions. No particular page or separately identified problem within this report is supplied.
\item Supplied literature-source link: Navarro--Ruhstorfer--Tiep--Vallejo, ``height-zero fields'': \url{https://doi.org/10.1016/j.aim.2024.109589}. This label is reproduced as supplied; no full bibliographic title or specific result connecting this reference to the four questions is provided.
\item Literature checked at: \texttt{2026-08-22}. This is supplied dated curation, not an independently verified assessment. The supplied status is ``open''. The supplied assessment states that these are broad unresolved questions concerning tensor products, vanishing sets, asymptotics, and modular character dimensions, and that current work resolves only restricted families rather than the general problem. The background explicitly calls this dated literature triage derived from supplied OWR curation and requests independent verification of time-sensitive status and cited claims.
\item The original statement describes these questions as having been hailed as definitive open problems in representation theory and algebraic combinatorics, and attributes deep connections with Lie theory, group theory, local-global conjectures in representation theory of finite groups, symplectic geometry, complexity theory, statistical mechanics, and quantum information theory. These are supplied contextual claims. The supplied Mathematics Subject Classification (2020) codes are \texttt{05E10}, \texttt{20C30}, and \texttt{20C08}.
\item The registry supplies the statement-status label \texttt{corrected\_verified}, with the qualification that only clear formula delimiters, notation, grammar, and typography were repaired, while mathematical content and scope were preserved. This records the supplied editorial verification claim; it does not establish a precise formulation or independent mathematical review.
\end{itemize}
\end{document}
